% TEMPLATE-MILC-v3.tex - plantilla maestra para todas las guias MILC v3
% Ver scripts/generadores/build-guia-milc.py para la lista completa de
% claves esperadas. El script valida que no queden placeholders sin
% reemplazar antes de escribir el archivo final.

\documentclass[11pt,letterpaper]{article}
\usepackage[margin=1.75cm]{geometry}
\usepackage[table]{xcolor}
\usepackage{fontspec}
\usepackage[spanish]{babel}
\usepackage{tikz}
\usepackage[most]{tcolorbox}
\usepackage{tabularx}
\usepackage{array}
\usepackage{fancyhdr}
\usepackage{titlesec}
\usepackage{ragged2e}
\usepackage{enumitem}
\usepackage{microtype}

\setmainfont{Helvetica Neue}
\setsansfont{Helvetica Neue}
\setlength{\parindent}{0pt}
\setlength{\parskip}{6pt}
\hyphenpenalty=200
\exhyphenpenalty=200
\pretolerance=400
\tolerance=2000
\setlength{\emergencystretch}{1em}
\renewcommand{\arraystretch}{1.25}
\setlist{leftmargin=5mm,itemsep=1mm,topsep=1mm}
\newcolumntype{Y}{>{\RaggedRight\arraybackslash}X}

\definecolor{milcMagenta}{HTML}{E6007E}
\definecolor{milcVino}{HTML}{5A0038}
\definecolor{milcTurquesa}{HTML}{0093A5}
\definecolor{milcVerde}{HTML}{58A923}
\definecolor{milcMostaza}{HTML}{E5B400}
\definecolor{milcOcre}{HTML}{B86600}
\definecolor{milcNegro}{HTML}{241816}
\definecolor{milcGris}{HTML}{F7F7F4}
\definecolor{milcLinea}{HTML}{E8E2DD}
\definecolor{milcGradePrimary}{HTML}{84CC16}
\definecolor{milcGradeDark}{HTML}{4D7C0F}
\definecolor{milcGradeSoft}{HTML}{ECFCCB}
\definecolor{milcGradeAccent}{HTML}{CFE7FF}
\definecolor{milcGradeText}{HTML}{FFFFFF}
\color{milcNegro}

\pagestyle{fancy}
\fancyhf{}
\lhead{\small Tecnología e Informática · Grado 7}
\rhead{\small Guía 18 · MILC v3}
\cfoot{\small\thepage}
\renewcommand{\headrulewidth}{0pt}

\newtcolorbox{titlebox}[1]{
  enhanced,colback=#1,colframe=#1,arc=7mm,boxrule=0pt,
  left=7mm,right=7mm,top=3mm,bottom=3mm,
  before skip=8pt,after skip=8pt,
  fontupper=\sffamily\bfseries\Large\color{white}
}

\newtcolorbox{softbox}[3]{
  enhanced,colback=#2,colframe=#1,borderline west={4pt}{0pt}{#1},
  arc=2mm,boxrule=.4pt,left=4mm,right=4mm,top=3mm,bottom=3mm,
  before skip=6pt,after skip=8pt,title={#3},coltitle=white,
  colbacktitle=#1,fonttitle=\sffamily\bfseries,fontupper=\RaggedRight,
  attach boxed title to top left={xshift=3mm,yshift=-2mm},
  boxed title style={arc=2mm, boxrule=0pt}
}

\newtcolorbox{drawbox}[2]{
  enhanced,colback=white,colframe=#1,arc=4mm,boxrule=1pt,
  left=5mm,right=5mm,top=4mm,bottom=4mm,before skip=8pt,after skip=8pt,
  title={#2},coltitle=white,colbacktitle=#1,fonttitle=\sffamily\bfseries,
  fontupper=\RaggedRight,
  attach boxed title to top left={xshift=4mm,yshift=-2mm},
  boxed title style={arc=3mm, boxrule=0pt}
}

\newtcolorbox{infoband}[1]{
  enhanced,colback=#1!8,colframe=#1!50,arc=2mm,boxrule=.5pt,
  left=4mm,right=4mm,top=2mm,bottom=2mm,before skip=4pt,after skip=4pt,
  fontupper=\small\RaggedRight
}

% Puente narrativo entre secciones (contrato editorial Conéctate).
% Da continuidad: "Acabas de X. Ahora vas a Y porque Z."
% Visualmente: banda izquierda en color del grado, texto en itálica pequeña.
\newtcolorbox{puentebox}{
  enhanced,colback=milcGradeSoft,colframe=milcGradeDark,
  borderline west={3pt}{0pt}{milcGradeDark},
  arc=0mm,boxrule=0pt,left=5mm,right=4mm,top=2.5mm,bottom=2.5mm,
  before skip=4pt,after skip=8pt,
  fontupper=\itshape\small\color{milcNegro}\RaggedRight
}
\newcommand{\puente}[1]{%
  \begin{puentebox}%
    \textcolor{milcGradeDark}{\textbf{\textrightarrow}}\hspace{2mm}#1%
  \end{puentebox}%
}

\newcommand{\linead}[1]{\noindent\color{milcLinea}\rule{#1}{0.5pt}\color{milcNegro}}
\newcommand{\respuesta}[1]{\par\vspace{2mm}\foreach \n in {1,...,#1}{\noindent\color{milcLinea}\rule{\linewidth}{0.5pt}\par\vspace{5mm}}\color{milcNegro}}
\newcommand{\checkbox}{\textcolor{milcGradeDark}{\fbox{\phantom{x}}}\hspace{5pt}}

\newcommand{\phasecard}[4]{
\begin{tcolorbox}[enhanced,colback=#3,colframe=#2,arc=3mm,boxrule=.5pt,height=3.05cm,valign=center,halign=center,left=1.5mm,right=1.5mm,top=2mm,bottom=2mm]
{\sffamily\bfseries\fontsize{22}{24}\selectfont\textcolor{#2}{#1}}\par
{\sffamily\bfseries\small #4}
\end{tcolorbox}
}

\begin{document}
\thispagestyle{empty}
\begin{tikzpicture}[remember picture,overlay]
  \fill[white] (current page.south west) rectangle (current page.north east);
  \path[fill=milcGradePrimary,rounded corners=16mm]
    ([xshift=.55cm,yshift=-.55cm]current page.north west) rectangle
    ([xshift=-.55cm,yshift=3.65cm]current page.south east);

  \node[anchor=north west,text=milcGradeText] at ([xshift=2.0cm,yshift=-1.85cm]current page.north west)
    {\sffamily\bfseries\fontsize{14}{16}\selectfont GRADO};
  \node[anchor=north west,text=milcGradeText,opacity=.82] at ([xshift=2.0cm,yshift=-2.75cm]current page.north west)
    {\sffamily\bfseries\fontsize{9}{11}\selectfont SERIE GUÍAS · TECNOLOGÍA E INFORMÁTICA};

  \begin{scope}[shift={([xshift=-3.05cm,yshift=-2.55cm]current page.north east)}]
    \fill[white,opacity=.80] (-.62,-.52) rectangle (.62,.52);
    \draw[milcGradeText,line width=1pt] (-.62,-.52) rectangle (.62,.52);
    \fill[milcGradeDark] (-.42,-.38) rectangle (-.22,.06);
    \fill[milcGradeAccent] (-.10,-.38) rectangle (.10,.24);
    \fill[milcGradePrimary!60!black] (.22,-.38) rectangle (.42,.38);
  \end{scope}

  \node[anchor=west,text=milcGradeText] at ([xshift=1.9cm,yshift=-12.85cm]current page.north west)
    {\sffamily\bfseries\fontsize{124}{124}\selectfont 7};
  \draw[milcGradeText,line width=1.7pt]
    ([xshift=4.52cm,yshift=-11.68cm]current page.north west) circle (.13cm);

  \node[anchor=west,text=milcGradeText,text width=8.7cm,align=left] at ([xshift=10.35cm,yshift=-11.6cm]current page.north west)
    {
     {\sffamily\bfseries\fontsize{19}{22}\selectfont Scratch ---\\mi primer\\programa\\con bloques}\\[4mm]
     {\sffamily\bfseries\fontsize{23}{26}\selectfont Guía 18-7-TIC}\\[2mm]
     {\sffamily\fontsize{13}{16}\selectfont Período 2 · Algoritmia y pensamiento computacional}\\[5mm]
     {\sffamily\bfseries\fontsize{11}{13}\selectfont Institución Educativa Sor María Juliana}\\[1mm]
     {\sffamily\fontsize{10}{12}\selectfont Autor: Dr. Álvaro Cárdenas Orozco}
    };

  \path[fill=milcGradeSoft,rounded corners=7mm]
    ([xshift=1.35cm,yshift=.85cm]current page.south west) rectangle
    ([xshift=-1.35cm,yshift=4.25cm]current page.south east);
  \draw[milcGradeDark,line width=.65pt,rounded corners=7mm]
    ([xshift=1.35cm,yshift=.85cm]current page.south west) rectangle
    ([xshift=-1.35cm,yshift=4.25cm]current page.south east);
  \node[anchor=center,text=milcNegro,text width=17.0cm,align=center] at ([yshift=2.55cm]current page.south)
    {\sffamily\fontsize{10}{12}\selectfont Metodología MILC · Apertura ancestral · Pensamiento computacional · Triángulo Dussel-Estoicismo-Floridi\\
    \textcolor{milcGradeDark}{\bfseries Producto:} Tu \textbf{primer programa en Scratch}: un \emph{sprite} (personaje) que: \textbf{(a) se mueve 100 pasos cuando se hace clic en bandera verde}, \textbf{(b) dice ``Hola'' por 2 segundos}, \textbf{(c) cambia de color y baila (repetir gira derecha-izquierda 10 veces)}. El programa guardado en tu cuenta de Scratch (\emph{scratch.mit.edu}) con tu cuenta o como invitado. Más \textbf{en el cuaderno}: captura impresa o boceto del proyecto + lista de los \textbf{6 bloques principales} que usaste con su categoría.};
\end{tikzpicture}
\mbox{}
\newpage

\section*{Apertura: Scratch --- mi primer programa con bloques}

\begin{softbox}{milcGradeDark}{milcGradeSoft}{Saber ancestral}
\textbf{Doña Mercedes la maestra rural de Cartago decía: ``Lo nuevo se aprende mejor jugando que estudiando.''} En la vereda La Plata, cuando enseñaba sumas a los niños del campo, no empezaba con ejercicios secos: \textbf{usaba piedrecitas, granos de maíz, tarjetas hechas a mano}. Los niños jugaban a sumar, restar, agrupar. \emph{Sin darse cuenta, estaban entendiendo matemáticas profundas}. La pedagogía del juego no era casualidad: doña Mercedes había observado que \emph{el cerebro aprende cuando se divierte}. Lo serio, presentado como serio, se aprende a medias; lo serio, presentado como juego, se aprende a fondo. \textbf{Esa intuición pedagógica es exactamente lo que motivó la creación de Scratch en el MIT}. Mitchel Resnick (su creador) sabía que enseñar programación a niños y jóvenes con código de texto sería frustrante. \emph{Diseñó Scratch para que sea juego primero, programación después}. Hoy abres ese taller-juego.
\end{softbox}

\begin{softbox}{milcTurquesa}{milcGris}{Saber contemporáneo}
\textbf{Scratch} es un lenguaje de programación visual desarrollado en el \emph{MIT Media Lab} (Massachusetts Institute of Technology) en 2007. Es \textbf{gratuito}, funciona en cualquier navegador (\texttt{scratch.mit.edu}), sin instalación. Más de \textbf{100 millones de proyectos} compartidos en su comunidad. Se basa en \textbf{bloques que se encajan como piezas de Lego}: cada bloque es una instrucción, los encajas en orden, formas un programa. \textbf{Conceptos visuales clave:} \emph{(1) Escenario (Stage)}: el fondo donde pasan las cosas. \emph{(2) Sprite}: un personaje u objeto (gato, perro, pelota). \emph{(3) Bloques}: instrucciones de colores. \emph{(4) Programa}: secuencia de bloques. \textbf{Categorías de bloques (9 colores):} Movimiento (azul), Apariencia (morado), Sonido (rosa), Eventos (amarillo), Control (naranja), Sensores (turquesa), Operadores (verde), Variables (rojo oscuro), Mis Bloques (rosa fuerte). Hoy usas bloques de las primeras 5 categorías.
\end{softbox}

\begin{softbox}{milcMostaza}{milcGris}{Pregunta-puente}
\textit{Cuando empieces a jugar con Scratch, ¿\textbf{cómo crees que se siente programar} con bloques en lugar de escribir código? ¿Más fácil? ¿Más difícil? ¿Más divertido?}
\end{softbox}

\begin{softbox}{milcVerde}{milcGris}{Saber y saber hacer}
Al terminar la clase: (1) podrás \textbf{identificar} la interfaz de Scratch (escenario, sprite, bloques); (2) sabrás \textbf{aplicar} los bloques de las 5 primeras categorías; (3) podrás \textbf{evaluar} si tu programa funciona correctamente; (4) habrás \textbf{creado} tu primer programa Scratch.
\end{softbox}

\small
\begin{tabularx}{\textwidth}{>{\bfseries}p{3.0cm}Y}
\rowcolor{milcGradePrimary!12}Autor & Dr. Álvaro Cárdenas Orozco\\
DBA & Construye algoritmos y diagramas de flujo para resolver problemas paso a paso (MEN, T\&I 7°)\\
Referentes & Saber ancestral del tejedor · Pensamiento computacional · Scratch\\
Producto final & Tu \textbf{primer programa en Scratch}: un \emph{sprite} (personaje) que: \textbf{(a) se mueve 100 pasos cuando se hace clic en bandera verde}, \textbf{(b) dice ``Hola'' por 2 segundos}, \textbf{(c) cambia de color y baila (repetir gira derecha-izquierda 10 veces)}. El programa guardado en tu cuenta de Scratch (\emph{scratch.mit.edu}) con tu cuenta o como invitado. Más \textbf{en el cuaderno}: captura impresa o boceto del proyecto + lista de los \textbf{6 bloques principales} que usaste con su categoría.\\
\end{tabularx}
\normalsize

\puente{Hoy haces 4 cosas: conoces la interfaz, exploras los bloques, creas tu primer programa, lo pruebas.}

\begin{titlebox}{milcGradeDark}
Ruta MILC: así vamos a aprender
\end{titlebox}
\begin{tabularx}{\textwidth}{>{\centering\arraybackslash}X>{\centering\arraybackslash}X>{\centering\arraybackslash}X>{\centering\arraybackslash}X}
\phasecard{1}{milcVerde}{milcGris}{Escuta\\[-1mm]\small Observo, escucho y pregunto.} &
\phasecard{2}{milcTurquesa}{milcGris}{Sistematización\\[-1mm]\small Pensamiento computacional.} &
\phasecard{3}{milcMagenta}{milcGris}{Praxis\\[-1mm]\small Construyo y aplico.} &
\phasecard{4}{milcVino}{milcGris}{Evaluación\\[-1mm]\small Reflexión liberadora.}
\end{tabularx}


\puente{Empezamos identificando partes de la interfaz de Scratch.}

\begin{titlebox}{milcVerde}
Fase 1 · Escuta
\end{titlebox}

\begin{softbox}{milcVerde}{milcGris}{Escena para escuchar}
\textbf{Actividad 1 · IDENTIFICA --- Interfaz de Scratch} (10 min · individual). Imagina (o el profe muestra) la pantalla de Scratch. En el cuaderno, dibuja un \textbf{rectángulo grande} representando la pantalla del computador. Adentro, marca con flechas y nombres las \textbf{4 zonas principales} de Scratch: \emph{(1) Escenario (Stage)}: el rectángulo grande arriba a la derecha, donde pasa la animación. \emph{(2) Sprites}: el panel debajo del escenario, donde se ven los personajes. \emph{(3) Categorías de bloques}: la columna de la izquierda con colores (azul, morado, rosa, etc.). \emph{(4) Área de programación}: el espacio central donde se arrastran los bloques para formar el programa.
\end{softbox}

\checkbox Dibujé un rectángulo representando la pantalla.\\
\checkbox Marqué las 4 zonas principales con flechas.\\
\checkbox Identifiqué para qué sirve cada zona.

\begin{infoband}{milcVerde}
\textbf{Lo que va al cuaderno (Actividad 1 · IDENTIFICA).} Título: «Actividad 1 --- Interfaz de Scratch». \textbf{Formato:} dibujo de la pantalla con 4 zonas marcadas. \textbf{Extensión:} 1/3 de página. \textbf{Sabes que terminaste cuando} sabes dónde está cada zona y qué hace.
\end{infoband}

\puente{Después aprendes las 5 categorías clave de bloques + cómo se encajan.}

\begin{titlebox}{milcTurquesa}
Fase 2 · Sistematización con pensamiento computacional
\end{titlebox}

\begin{softbox}{milcTurquesa}{milcGris}{Cuatro pilares aplicados al tema}
Ahora aprende los \textbf{bloques de las 5 categorías} que vas a usar hoy + cómo se encajan.
\end{softbox}

\small
\begin{tabularx}{\textwidth}{>{\bfseries\RaggedRight\arraybackslash}p{3.6cm}Y}
\rowcolor{milcTurquesa!90}\color{white}Pilar & \color{white}\bfseries Aplicación al tema\\
1. Descomposición & \textbf{Categoría 1 --- Eventos (amarillo).} \textbf{El programa empieza con un evento.} Bloque clave: \emph{``Al hacer clic en la bandera verde''} (la bandera arriba a la izquierda del escenario). Cuando el usuario hace clic en esa bandera, los bloques que están debajo se ejecutan. Otros eventos: \emph{``Al presionar tecla [letra]''}, \emph{``Al hacer clic en este sprite''}. Los eventos son como el botón START de un juego.\\
2. Reconocimiento de patrones & \textbf{Categoría 2 --- Movimiento (azul).} \textbf{Mueven al sprite por el escenario.} Bloques clave: \emph{``Mover 10 pasos''} (avanza), \emph{``Girar 15 grados a la derecha''} (rota), \emph{``Ir a X, Y''} (teletransporta a una coordenada). El escenario es un plano cartesiano: el centro es (0,0), derecha es X positivo, arriba es Y positivo. \textbf{Categoría 3 --- Apariencia (morado).} \textbf{Cambian cómo se ve el sprite.} Bloques clave: \emph{``Decir 'Hola' por 2 segundos''} (globo de texto), \emph{``Cambiar disfraz''} (cambiar imagen), \emph{``Cambiar tamaño''} (más grande o pequeño), \emph{``Cambiar color''}.\\
3. Abstracción & \textbf{Categoría 4 --- Sonido (rosa).} \textbf{Reproducen sonidos.} Bloques clave: \emph{``Tocar sonido [nombre]''}, \emph{``Reproducir sonido hasta acabar''}, \emph{``Detener sonido''}. Scratch incluye una biblioteca de sonidos prefabricados; también puedes grabar tus propios sonidos con el micrófono. \textbf{Categoría 5 --- Control (naranja).} \textbf{Estructuras de control: condicionales, bucles, esperas.} Bloques clave: \emph{``Esperar 1 segundo''} (pausa), \emph{``Repetir 10 veces''} (bucle FOR de la S7), \emph{``Repetir hasta que [condición]''} (bucle WHILE de la S7), \emph{``Si [condición] entonces [acción] si no [acción]''} (condicional de la S6). \textbf{Las 5 categorías más usadas para tu primer programa son estas.}\\
4. Algoritmo conceptual & \textbf{Cómo se encajan los bloques.} Los bloques tienen formas que indican cómo se conectan: \textbf{(1) Sombreros (eventos):} forma redondeada arriba, recta abajo. Solo van al principio del programa. \textbf{(2) Bloques sándwich (acciones):} con muescas arriba y abajo, se encajan en cualquier punto del medio. \textbf{(3) Bloques con espacio interno (bucles, condicionales):} tienen un hueco grande adentro donde se meten otros bloques. \emph{Arrastras un bloque y lo encajas con un sonido de clic}. Si dos bloques no encajan visualmente, no se conectan.\\
\end{tabularx}
\normalsize

\begin{drawbox}{milcTurquesa}{5 categorías + encajes}
\textbf{Eventos (amarillo):} Al hacer clic en bandera verde, al presionar tecla. \\[1mm]
\textbf{Movimiento (azul):} mover, girar, ir a X-Y. \\[1mm]
\textbf{Apariencia (morado):} decir, cambiar disfraz, color, tamaño. \\[1mm]
\textbf{Sonido (rosa):} tocar sonido. \\[1mm]
\textbf{Control (naranja):} esperar, repetir, si...entonces. \\[1mm]
\textbf{Encajes:} sombreros (arriba) · sándwich (medio) · con hueco interno (control).
\end{drawbox}

\begin{softbox}{milcMostaza}{milcGris}{Errores comunes y su corrección}
\textbf{Errores frecuentes al empezar con Scratch.} (1) \emph{Olvidar el bloque de evento (bandera verde)} --- el programa no arranca. Siempre empieza con un bloque amarillo. (2) \emph{Encajar bloques sin que hagan clic} --- los bloques no se conectaron, el segundo no se ejecutará. (3) \emph{Querer hacer un videojuego en la primera sesión} --- empieza con un sprite que se mueve y dice algo. La complejidad viene en S9 y S10. (4) \emph{No probar el programa} --- haces clic en la bandera verde para ver qué pasa. Si no se ejecuta como esperabas, revisas y corriges. (5) \emph{Olvidar guardar} --- en Scratch en línea se guarda automático con cuenta. Sin cuenta, no se guarda al cerrar.
\end{softbox}

\begin{infoband}{milcTurquesa}
\textbf{Lo que va al cuaderno (Actividad 2 · EXPLICA).} Título: «Actividad 2 --- 5 categorías de Scratch». \textbf{Formato:} tabla de 5 filas, 3 columnas (categoría, color, 2-3 bloques principales). \textbf{Extensión:} 5 filas. \textbf{Sabes que terminaste cuando} reconoces cualquier bloque y sabes a qué categoría pertenece por su color.
\end{infoband}

\puente{Con todo claro, abres Scratch y haces tu primer programa.}

\begin{titlebox}{milcMagenta}
Fase 3 · Praxis con pensamiento computacional
\end{titlebox}

\begin{softbox}{milcMagenta}{milcGris}{Construyendo el producto con los cuatro pilares}
\textbf{Actividad 3 · CREA --- Tu primer programa Scratch} (35 min · individual). Vas a abrir Scratch y crear tu primer programa real. Es \textbf{el momento esperado del periodo}: vas a programar de verdad.
\end{softbox}

\small
\begin{tabularx}{\textwidth}{>{\bfseries\RaggedRight\arraybackslash}p{3.6cm}Y}
\rowcolor{milcMagenta!90}\color{white}Pilar & \color{white}\bfseries Aplicación al producto\\
Descomposición & \textbf{Paso 1 --- Abre Scratch.} En el navegador escribe \texttt{scratch.mit.edu}. Haz clic en \emph{``Crear''} (esquina superior). Aparece la interfaz: escenario arriba a la derecha con el gato amarillo, panel de bloques a la izquierda, área de programación en el centro. Si tienes cuenta, inicia sesión (para que se guarde automáticamente). Si no, sigue como invitado.\\
Patrones & \textbf{Paso 2 --- Familiarízate (2-3 min).} Arrastra cualquier bloque del panel de la izquierda al área de programación central. Haz clic en él (se ejecuta una vez). Trata de encajar otro bloque debajo. Cuando dos bloques se encajan correctamente, se ve un clic visual. Esto te entrena en la mecánica básica.\\
Abstracción & \textbf{Paso 3 --- Programa 1: el sprite se mueve y saluda.} En el área de programación arma este programa de \textbf{4 bloques}: \emph{(a) ``Al hacer clic en bandera verde'' [Eventos, amarillo].} \emph{(b) ``Decir 'Hola' por 2 segundos'' [Apariencia, morado].} \emph{(c) ``Mover 100 pasos'' [Movimiento, azul].} \emph{(d) ``Decir 'Listo!' por 2 segundos'' [Apariencia, morado].}. Encájalos uno debajo del otro. \textbf{Prueba:} haz clic en la bandera verde (arriba del escenario, ícono de bandera). El gato dirá Hola, se moverá 100 pasos a la derecha, dirá Listo.\\
Algoritmo de armado & \textbf{Paso 4 --- Programa 2: agregar el baile.} Debajo del bloque \emph{``Mover 100 pasos''}, agrega: \emph{(e) ``Cambiar efecto color por 25'' [Apariencia, morado].} \emph{(f) ``Repetir 10 veces'' [Control, naranja]}. Dentro del bloque de repetir, mete: \emph{(g) ``Girar 15 grados a la derecha'' [Movimiento, azul].} \emph{(h) ``Esperar 0.1 segundos'' [Control, naranja].} \emph{(i) ``Girar 15 grados a la izquierda'' [Movimiento, azul].} Cuando los pongas dentro del bucle ``Repetir 10 veces'', el bucle los abrazará. Prueba de nuevo: el sprite ahora hace el baile completo.\\
\end{tabularx}
\normalsize

\begin{drawbox}{milcMagenta}{Antes de cerrar la S8}
\checkbox Abrí Scratch en scratch.mit.edu. \\[1mm]
\checkbox Mi programa tiene bloque de evento (bandera verde). \\[1mm]
\checkbox El sprite se mueve, dice Hola y baila. \\[1mm]
\checkbox El bucle de repetir 10 veces abraza los bloques de giro. \\[1mm]
\checkbox Personalicé el sprite y/o el fondo. \\[1mm]
\checkbox Captura o boceto del programa en el cuaderno.
\end{drawbox}

\begin{softbox}{milcMagenta}{milcGris}{Plantilla de trabajo}
\textbf{Estructura del programa.} \textit{Bloque 1 (Eventos):} Al hacer clic en bandera verde. \textit{Bloque 2 (Apariencia):} Decir Hola 2 seg. \textit{Bloque 3 (Movimiento):} Mover 100 pasos. \textit{Bloque 4 (Apariencia):} Decir Listo 2 seg. \textit{Bloque 5 (Apariencia):} Cambiar color +25. \textit{Bloque 6 (Control - bucle):} Repetir 10 veces [Girar 15° derecha, Esperar 0.1s, Girar 15° izquierda].
\end{softbox}

\begin{infoband}{milcMagenta}
\textbf{Lo que va al cuaderno (Actividad 3 · CREA).} Título: «Actividad 3 --- Mi primer programa Scratch». \textbf{Formato:} captura + lista de 6 bloques con categoría/color + reflexión. \textbf{Extensión:} 1 página. \textbf{Sabes que terminaste cuando} tu programa se ejecuta correctamente al hacer clic en bandera verde.
\end{infoband}

\puente{El producto es el programa guardado + captura + lista de bloques en cuaderno.}

\begin{titlebox}{milcMostaza}
Producto de la sesión
\end{titlebox}
\textbf{Mi primer programa Scratch · sprite que se mueve, saluda y baila}.
\begin{softbox}{milcMostaza}{milcGris}{Criterios de calidad}
(1) El programa tiene bloque de evento (bandera verde). (2) El sprite ejecuta las 3 acciones (moverse, saludar, bailar). (3) El bucle de repetir 10 veces abraza correctamente los bloques de giro. (4) Se personalizó el sprite y/o fondo. (5) En el cuaderno hay captura + lista de 6 bloques + reflexión firmada.
\end{softbox}

\puente{Antes de salir, ejecuta el programa al menos 3 veces para verificar que funciona.}

\begin{titlebox}{milcVino}
Fase 4 · Evaluación liberadora — cinco dimensiones
\end{titlebox}

\begin{softbox}{milcVino}{milcGris}{Sentido de la evaluación}
Evaluar no es solo revisar una nota. Es reconocer qué aprendí, qué cambió en mí, qué emociones manejé, cómo me relaciono con la ciudad y la comunidad, y qué le dejaré a quien viene detrás.
\end{softbox}

\textbf{Mi reflexión liberadora en cinco dimensiones.} Completa cada frase iniciada:

\small
\begin{tabularx}{\textwidth}{>{\bfseries\RaggedRight\arraybackslash}p{3.4cm}Y>{\RaggedRight\arraybackslash}p{4.6cm}}
\rowcolor{milcVino}\color{white}Dimensión & \color{white}\bfseries Mi respuesta (frame starter) & \color{white}\bfseries Algo que descubrí\\
1. Personal & Lo que más cambió en mí fue \linead{6.5cm} & \linead{4.2cm}\\
2. Emocional & La emoción más fuerte fue \linead{2.5cm} y la regulé \linead{3cm} & \linead{4.2cm}\\
3. Ciudadana & En democracia esto importa porque \linead{6cm} & \linead{4.2cm}\\
4. Local (Cartago) & En mi barrio sería útil para \linead{6.5cm} & \linead{4.2cm}\\
5. Intergeneracional & A mi abuela/abuelo le diría \linead{6.5cm} & \linead{4.2cm}\\
\end{tabularx}
\normalsize

\begin{infoband}{milcVino}
\textbf{Para la próxima sustentación.} Vuelve a esta tabla en seis meses. Verás que las respuestas cambian — no porque lo aprendido cambie, sino porque tú cambias. La evaluación liberadora es una conversación contigo a través del tiempo.
\end{infoband}


\puente{Cerramos con 3 voces sobre por qué programar es una de las habilidades clave del siglo XXI.}

\begin{titlebox}{milcGradeDark}
Triángulo de pensamiento · cierre filosófico
\end{titlebox}

\begin{softbox}{milcVino}{milcGris}{Dussel · lente del nosotros}
\textit{````Scratch democratiza la programación: lo que era para ingenieros con doctorado ahora es para niños del barrio. Eso es justicia educativa.''''}

Antes de Scratch, \textbf{programar requería años de estudio formal}: aprender sintaxis compleja, depurar errores de texto, leer documentación en inglés. \emph{Eso excluía a niños, a estudiantes de barrios populares, a personas sin acceso a universidades caras}. Scratch puso la programación al alcance de \textbf{cualquier niño con computador y conexión}. La abuela del Cauca podía aprender; tú lo haces hoy. \emph{El MIT, una de las universidades más exclusivas del mundo, creó Scratch como acto de justicia educativa}.

\textbf{¿Estoy aprovechando este acceso al saber que generaciones anteriores no tuvieron?}
\end{softbox}
\linead{\linewidth}\\[1mm]
\linead{\linewidth}

\begin{softbox}{milcMostaza}{milcGris}{Estoicismo · Marco Aurelio (emperador de los hábitos diarios) · cuidado interior}
\textit{````El que practica una habilidad 30 minutos al día durante 1 año, llega lejos. Scratch te da las herramientas para esa práctica.''''}

\textbf{Scratch es la herramienta. La práctica es tuya}. Marco Aurelio escribía 30 minutos al día durante décadas. Tú puedes practicar Scratch 20-30 minutos al día. En 1 año, programas con fluidez. En 3 años, eres capaz de hacer cosas que adultos no saben hacer. \emph{La disciplina pequeña diaria construye habilidades enormes a largo plazo}.

\textbf{¿Qué pasaría si practicara Scratch 20 minutos cada día en lugar de solo en clase?}
\end{softbox}
\linead{\linewidth}\\[1mm]
\linead{\linewidth}

\begin{softbox}{milcTurquesa}{milcGris}{Floridi · lente de la infoesfera}
\textit{````En 2030, no saber programar mínimamente será como en 1990 no saber leer. Scratch es el lápiz para esa nueva alfabetización.''''}

Las generaciones de tus abuelos aprendieron a leer; las de tus papás aprendieron a usar computadoras; \textbf{tu generación aprende a programar}. No para ser programador profesional, sino para \emph{entender el mundo que te rodea}. Scratch es el primer paso. En grados 8°-11° vas a aprender más lenguajes. En la universidad, código real. \emph{Pero el oficio fundamental --- pensar en bloques de instrucciones --- lo aprendes hoy}.

\textbf{¿Estoy reconociendo lo importante que es esta alfabetización para mi vida adulta?}
\end{softbox}
\linead{\linewidth}\\[1mm]
\linead{\linewidth}

\begin{titlebox}{milcVino}
Compromiso final
\end{titlebox}

\small
\begin{tabularx}{\textwidth}{>{\RaggedRight\arraybackslash}p{3.0cm}YYY}
\rowcolor{milcVino}\color{white}\bfseries Criterio & \color{white}\bfseries Lo logré & \color{white}\bfseries En proceso & \color{white}\bfseries Necesito apoyo\\
Comprensión del tema & Explico ideas centrales con mis palabras. & Reconozco ideas pero falta precisión. & Necesito repasar conceptos base.\\
Pensamiento computacional & Apliqué los 4 pilares al diseñar mi producto. & Apliqué algunos pilares. & Necesito releer la fase 2.\\
Aplicación y producto & Mi producto cumple los criterios y se puede revisar. & Producto funcional con ajustes pendientes. & Aún en construcción.\\
Reflexión 5 dimensiones & Respondí cada dimensión con honestidad. & Respondí algunas con detalle. & Mis respuestas son muy generales.\\
Triángulo de pensamiento & Conecté las 3 voces con mi trabajo real. & Conecté al menos una. & No relaciono las voces con mi proceso.\\
\end{tabularx}
\normalsize

\textbf{Hoy entendí que:}\\[1mm]
\linead{\linewidth}\\[1mm]
\linead{\linewidth}

\textbf{Mi compromiso para la próxima sesión es:}\\[1mm]
\linead{\linewidth}\\[1mm]
\linead{\linewidth}

\begin{softbox}{milcMostaza}{milcGris}{Cita de despedida · Marco Aurelio}
\textit{``El obstáculo en el camino se vuelve el camino. Lo que estorba al camino se vuelve camino.''}\\[1mm]
Cada error de hoy, bien observado, se convierte en pieza del camino que pisarás mañana.
\end{softbox}

\small
\begin{tabularx}{\textwidth}{>{\bfseries}p{3.6cm}Y}
\rowcolor{milcGradeDark!12}Nombre completo: & \linead{10cm}\\
Fecha: & \linead{4cm} \quad Curso/grupo: \linead{4cm}\\
Mi firma: & \linead{9cm}\\
\end{tabularx}
\normalsize

\begin{infoband}{milcGradeDark}
\textbf{Para la próxima sesión.} Esta semana, abre Scratch otras 2 veces. Modifica tu programa: cambia los números, agrega más bloques, prueba qué pasa. La curiosidad es tu motor. La próxima clase aprendes a hacer una \textbf{historia interactiva}: cuento donde el usuario escoge qué pasa.
\end{infoband}

\end{document}
